% PILOTO INCREMENTAL, basado directamente en temas/tema1.tex.
% Conserva la portada, las 21 diapositivas originales, sus figuras y ejercicios.
% Añade nueve diapositivas: índice, puentes, ejemplos y práctica graduada.
% Los marcadores "Diapositiva original N" garantizan trazabilidad.
% Archivo independiente: no se modifica el temario oficial ni otros pilotos.
% Compilar desde la raíz del repositorio.
\documentclass[main.tex]{subfiles}

\begin{document}
\graphicspath{{imagenes/tema1/}{../imagenes/tema1/}}

% Diapositiva original 1
\portadatema{1}{Funciones de varias variables}
% NUEVA: Índice general
\begin{frame}{Contenido del tema}
\begin{enumerate}
\item Funciones reales de varias variables.
\item Dominio, imagen y representación gráfica.
\item Determinación del dominio: ejemplos.
\item Operaciones con funciones.
\item Curvas y conjuntos de nivel.
\item Ejercicios de aplicación.
\end{enumerate}
\medskip
\small Los conceptos se ilustrarán también con aplicaciones de visualización científica e informática.
\end{frame}

% Diapositiva original 2
\begin{frame}{Funciones de varias variables}
La temperatura $T$ en un punto de la superficie de la Tierra a una hora determinada depende de la longitud y la latitud de dicho punto.
\medskip

Se puede pensar en $T$ como una función de dos variables y, si las denominamos $x$ e $y$, esta dependencia funcional se puede poner como
\[T=f(x,y).\]
\end{frame}
% Diapositiva original 3
\begin{frame}{Funciones de varias variables}
El volumen de una caja rectangular de dimensiones $x,y,z$ es $\text{Volumen}=x\cdot y\cdot z$.
\medskip

Este es un ejemplo de una función real de tres variables reales, que simbolizamos por:
\[f(x,y,z)=x\cdot y\cdot z.\]
\figura{s3_2.png}{.45}{.43}
\end{frame}
% Diapositiva original 4
\begin{frame}{Funciones de varias variables}
\textbf{Definición}
\medskip

En general, una función real de $n$ variables reales es una correspondencia que asigna a cada $n$-tupla $(x_1,x_2,\ldots,x_n)$ de su dominio un único valor
\[u=f(x_1,x_2,\ldots,x_n).\]
\end{frame}
% Diapositiva original 5
\begin{frame}{De una variable a varias variables}
Los conceptos conocidos para funciones de una variable tienen su equivalente para funciones de $n$ variables.
\medskip

Las entradas pasan de ser números $x$ a vectores
$(x_1,\ldots,x_n)$, pero conservamos las ideas de
\textbf{dominio}, \textbf{imagen} y \textbf{gráfica}.
\medskip

Por ejemplo, para $f(x,y)=x^2+y^2$:
\[
D=\mathbb R^2,\qquad \operatorname{Im}(f)=[0,+\infty).
\]
El codominio puede ser $\mathbb R$, aunque la imagen sea menor.
\end{frame}
% Diapositiva original 6
\begin{frame}{Funciones de varias variables}
Particularizando a dos variables, una función $f$ de dos variables asigna a cada par ordenado de números reales $(x,y)$ perteneciente a un conjunto $D$ un número real único denotado por $u=f(x,y)$.
\medskip

El conjunto $D$ es el dominio de $f$ y su imagen es el conjunto de valores que toma $f$.
\figura{s6_0_v2.jpg}{.90}{.30}
\end{frame}
% Diapositiva original 7
\begin{frame}{Funciones de varias variables}
Las funciones de dos variables se suelen escribir de la forma $z=f(x,y)$ para explicitar el valor tomado por la función $f$ en un punto $(x,y)$.
\medskip

Las variables $x$ e $y$ son las variables independientes y $z$ es la variable dependiente.
\end{frame}
% NUEVA: Ejemplo de base antes de las gráficas
\begin{frame}{Ejemplo: evaluación y dominio}
Consideremos la función $f(x,y)=x^2+2y$.
\medskip

Evaluamos sustituyendo las dos variables:
\[
f(1,2)=1^2+2\cdot2=5,\qquad
f(-1,3)=(-1)^2+2\cdot3=7.
\]
\medskip

Su \textbf{dominio natural} es todo $\mathbb R^2$, porque no aparecen operaciones que impongan restricciones.
\medskip

Este ejemplo servirá como referencia al estudiar casos más complejos.
\end{frame}

% Diapositiva original 8
\begin{frame}{Representación gráfica de funciones de dos variables}
\small Si $f$ es una función de dos variables con dominio $D$, entonces la gráfica de $f$ es el conjunto de ternas $(x,y,z)$ de $\mathbb{R}^3$ tales que $z=f(x,y)$ y $(x,y)\in D$.
\medskip

Igual que las gráficas de funciones de una variable son curvas en el plano, las gráficas de funciones de dos variables son superficies en el espacio.
\figura{s8_2.png}{.44}{.47}
\end{frame}
% Diapositiva original 9
\begin{frame}{Ejemplo: superficie en el espacio}
Por ejemplo, $z=f(x,y)=x+y^2$ es una función de dos variables con dominio todos los pares de números reales y cuya representación en 3D es
\figura{s9_2.png}{.85}{.58}
\end{frame}
% NUEVA: Contexto para los ejemplos originales de dominio
\begin{frame}{Determinación del dominio natural}
El ejemplo polinómico anterior tiene dominio $\mathbb R^2$.
En otras expresiones, las operaciones introducen restricciones.
\medskip
\begin{itemize}
\item En un \textbf{cociente}, el denominador no puede ser cero.
\item En una \textbf{raíz cuadrada}, el radicando debe ser no negativo.
\item En un \textbf{logaritmo}, el argumento debe ser positivo.
\end{itemize}
\medskip

Si aparecen varias condiciones, el dominio es su
\textbf{intersección}. Veamos tres ejemplos.
\end{frame}

% Diapositiva original 10
\begin{frame}{Ejemplo: dominio de una función racional}
\textbf{Ejemplo}
\[z=f(x,y)=\frac{3x-5y}{y-x^2},\quad D=\{(x,y)\in\mathbb{R}^2:y-x^2\ne0\}.\]
\figura{s10_4.png}{.70}{.44}
El dominio es todo el plano $\mathbb{R}^2$ excepto los puntos de la parábola $y=x^2$.
\end{frame}
% Diapositiva original 11
\begin{frame}{Ejemplo: dominio de una función radical}
\textbf{Ejemplo}
\[z=f(x,y)=\sqrt{9-x^2-y^2},\quad D=\{(x,y)\in\mathbb{R}^2:9-x^2-y^2\ge0\}.\]
\figura{s11_5.png}{.77}{.40}
\small Como $x^2+y^2=9$ es la ecuación de una circunferencia de centro el origen y radio $3$, $D$ representa el interior y la frontera de dicha circunferencia.
\end{frame}
% Diapositiva original 12
\begin{frame}{Ejemplo: dominio de una función de tres variables}
El dominio de la función
\[f(x,y,z)=\ln(z-y)+xy\sin z\]
es
\[D=\{(x,y,z)\in\mathbb{R}^3:z>y\}.\]
Es decir, todos los puntos que están por encima del plano $z=y$.
\end{frame}
% NUEVA: Dominio con restricciones combinadas
\begin{frame}{Ejemplo adicional: dos restricciones simultáneas}
Consideremos
\[
f(x,y)=\frac{\sqrt{4-x^2-y^2}}{y-1}.
\]
Debe cumplirse simultáneamente:
\[
4-x^2-y^2\geq0,\qquad y-1\ne0.
\]
Por tanto,
\[
\boxed{D=\{(x,y)\in\mathbb R^2:
x^2+y^2\leq4,\ y\ne1\}.}
\]
\medskip

Es el disco cerrado de radio $2$, del que se excluyen los puntos de la recta $y=1$ que pertenecen al disco.
\end{frame}

% Diapositiva original 13
\begin{frame}{Operaciones con funciones de varias variables}
Las funciones de varias variables pueden operarse de igual forma que las de una variable, ya sea con suma, producto, cociente o composición de funciones.
\medskip

Para $f(x,y)=x+y$ y $g(x,y)=x^2+y^2$:
\[
(f+g)(x,y)=x+y+x^2+y^2,
\qquad
\left(\frac fg\right)(x,y)=\frac{x+y}{x^2+y^2}.
\]
La suma está definida en $\mathbb R^2$;
en el cociente hay que excluir $(0,0)$, donde $g=0$.
\end{frame}
% NUEVA: Puente conceptual hacia las curvas de nivel
\begin{frame}{De la superficie a la representación plana}
Las gráficas de funciones de dos variables se representan generalmente mediante superficies tridimensionales.
\medskip

En cartografía, representación de datos y visualización científica
interesa también describirlas mediante dibujos en el plano.
\medskip

Una técnica consiste en representar los puntos del dominio donde la función toma \textbf{el mismo valor constante}.
\medskip

Esta idea conduce a las \textbf{curvas de nivel}.
\end{frame}

% Diapositiva original 14
\begin{frame}{Curvas de nivel: definición e interpretación}
Las curvas de nivel de una función $f$ de dos variables son las curvas con ecuaciones $f(x,y)=k$, donde $k$ es una constante que pertenece a la imagen de $f$.
\begin{columns}[c]
\begin{column}{.48\textwidth}\figura{s14_3.png}{.85}{.56}\end{column}
\begin{column}{.48\textwidth}\figura{s14_5.png}{.98}{.56}\end{column}
\end{columns}
\end{frame}
% NUEVA: Primer ejemplo de curvas de nivel sencillo
\begin{frame}{Ejemplo adicional: niveles de una función afín}
Sea $f(x,y)=x+y$. Para cada nivel $k$:
\[
f(x,y)=k\quad\Longleftrightarrow\quad y=k-x.
\]
Las curvas de nivel son \textbf{rectas paralelas}.
\[
\begin{array}{c|c}
k & \text{curva de nivel}\\\hline
-1 & y=-1-x\\
0 & y=-x\\
1 & y=1-x
\end{array}
\]
\medskip

El siguiente ejemplo original permite estudiar una geometría diferente: las circunferencias.
\end{frame}

% Diapositiva original 15
\begin{frame}{Ejemplo: curvas de nivel de una semiesfera}
\textbf{Ejemplo}
\vfill
Obtener las curvas de nivel de
\[z=\sqrt{100-x^2-y^2}.\]
\vfill
\end{frame}
% Diapositiva original 16
\begin{frame}{Resolución: niveles de la semiesfera}
\small\textbf{Ejemplo}\quad $z=\sqrt{100-x^2-y^2}$.
\medskip

Sustituimos la variable $z$ por una constante $z=c$:
\[c=\sqrt{100-x^2-y^2}.\]
Elevando al cuadrado,
\[c^2=100-x^2-y^2.\]
Reorganizando,
\[x^2+y^2=100-c^2.\]
Se obtienen circunferencias de centro el origen para $0\le c<10$. Para $c=10$ se obtiene el punto $(0,0)$. Para $c<0$ o $c>10$ no hay curvas de nivel.
\end{frame}
% Diapositiva original 17
\begin{frame}{Representación gráfica de los niveles}
\textbf{Interpretación gráfica}\par
Representar las curvas de nivel de la función de dos variables
\begin{columns}[c]
\begin{column}{.43\textwidth}
\[z=\sqrt{100-x^2-y^2}\]
\[x^2+y^2=100-c^2\]
\end{column}
\begin{column}{.55\textwidth}\figura{s17_3.png}{1}{.58}\end{column}
\end{columns}
\end{frame}
% Diapositiva original 18
\begin{frame}{Relación entre superficie y curvas de nivel}
\figura{s18_0.jpg}{.78}{.83}
\end{frame}
% Diapositiva original 19
\begin{frame}{Curvas de nivel de una función no radial}
\figura{s19_1.jpg}{1}{.67}\begin{center}\scriptsize Curvas y superficies de $f(x,y)=-xy\,e^{-x^2-y^2}$.\end{center}
\end{frame}
% Diapositiva original 20
\begin{frame}{Curvas de nivel de una función racional}
\figura{s20_1.jpg}{1}{.76}
\end{frame}
% Diapositiva original 21
\begin{frame}{Aplicación: curvas de nivel en cartografía}
\figura{s21_0.jpg}{.73}{.70}
\begin{center}\scriptsize Las curvas unen puntos de igual altitud; su proximidad informa sobre la pendiente.\end{center
\end{frame}
% NUEVA: Informática sin cambiar el hilo matemático
\begin{frame}{Aplicación informática: funciones de coste}
En optimización y aprendizaje automático, una función de coste depende de los parámetros de un modelo.
\medskip

Por ejemplo:
\[
L(w_1,w_2)=(w_1-2)^2+2(w_2+1)^2.
\]
Sus curvas de nivel verifican
\[
(w_1-2)^2+2(w_2+1)^2=k.
\]
\medskip

Para $k>0$ son \textbf{elipses} centradas en $(2,-1)$.
Estas curvas permiten interpretar cómo varía el coste al modificar los parámetros.
\end{frame}

% NUEVA: Práctica complementaria básica
\begin{frame}{Ejercicios complementarios: nivel básico}
\begin{enumerate}
\item Si $f(x,y)=2x+y$, calcular $f(1,2)$ y $f(-1,3)$.
\medskip
\item Determinar el dominio natural de
\[
g(x,y)=\sqrt{9-x^2-y^2}.
\]
\item Describir las curvas de nivel $k=0,1$ de
\[
h(x,y)=x+y.
\]
\end{enumerate}
\medskip
\small Estos ejercicios trabajan directamente los conceptos y procedimientos elementales.
\end{frame}

% NUEVA: Práctica complementaria intermedia
\begin{frame}{Ejercicios complementarios: nivel intermedio}
\begin{enumerate}
\item Determinar y dibujar el dominio de $f(x,y)=\ln(y-x^2)$.
\medskip
\item Encontrar el dominio de
\[
g(x,y)=\frac{\sqrt{4-x^2-y^2}}{y-1}.
\]
\item Describir geométricamente los niveles $k=0,1,4$ de
\[
h(x,y)=x^2+y^2.
\]
\end{enumerate}
\end{frame}

% Diapositiva original 22
\begin{frame}{Ejercicios originales: consolidación y ampliación}
\small\textbf{Ejercicios}
\medskip

Encontrar el dominio de la función
\[f(x,y,z)=\arcsin(xy)\,e^{3z}+e^{(y+z)\ln(x+z)}.\]
Describir los conjuntos de nivel para los valores de $k$ indicados:
\begin{enumerate}[a)]
\item $f(x,y)=x^2+y^2$,\quad $k=0,1,2,3$.
\item $f(x,y,z)=x^2+y^2$,\quad $k=0,1,2$.
\item $f(x,y)=x^2-y^2$,\quad $k=-2,-1,0,1,2$.
\end{enumerate}
En a) y c) se trata de curvas de nivel; en b), de superficies de nivel.
\end{frame}
\end{document}
